% !TEX root = main.tex
\section{적용 범위와 남은 검증 과제}

\subsection{디지털 회계와 실제 제품의 경계}
zkDPP의 검증 대상은 원장이 수락한 디지털 전이와 제출 값의
관계이다. 초기 등록과 실제 측정의 진실성, 공정 수행, 재료 적
합성, 물리적 제품과 문서의 연결은 외부 신뢰에 의존한다. 자
산 역할은 후속 연산의 허용 범위를 구분하며 물질 종류를 증명
하지 않는다. 가공이 문서 해시를 바꿀 수 있다는 사실도 물질
의 물리적 전환을 증명하는 것은 아니다.

탄소는 수용한 초기 값과 이후 운송·공정 증분을 배분한 회계
값이다. 초기 등록의 기본값 0은 등록 전 배출이 없었다는 증거
가 아니고, 명세가 모든 생애주기 배출 항목을 자동 포함하는 것
도 아니다. 폐기물의 \(a=e=0\)과 종료 전용 소비는 참조 회계
규칙이다. 이를 실제 폐기물의 환경 영향이 없다는 판단으로 확
장하지 않는다. DPP 검증 역시 선택된 세 필드와 제품 주장을
결합하며, 문서 전체의 정확성이나 규제·인증 체계 전체의 준수
를 보장하지 않는다.

디지털 회계 내부에서도 개별 전이의 합계 보존과 최종 제품
군의 탄소 귀속은 구분해야 한다. 식 (19)에서 첫 일반용 출력
의 질량이 0이어도 반올림 잔여를 받기 때문이다. 예를 들어
승인된 \((m,n)=(1,3)\) 정책에서 입력이 \((q,a,e)=(2,2,1)\),
질량 손실과 공정 증분이 0이고 세 출력이 모두 일반용이라
고 하자. 출력 질량 \((0,1,1)\)과 크레딧 \((0,1,1)\)에 대응하는
탄소는 \((1,0,0)\)이다. 각 전이 관계는 만족하지만 첫 출력을
종료하면 양의 질량 출력의 탄소 합은 0이 되며, 그 출력은
\(τ_{R}=1,τ_{C}=0\)인 발급 조건도 만족한다. 이는 물리 측정의 오
류가 아니라 참조 정수 배분 규칙이 허용하는 실행이다. 따라서
본 구성의 발급은 현재 객체의 승인된 정책 조건을 의미하며, 경
로에 불변인 탄소 집약도나 배분 조작 저항성을 보장하지 않는
다. 이 사례를 금지하려면 질량 0 출력과 잔여 배분에 관한 추
가 정책이 필요하다. 본 논문은 기존 규칙을 바꾸어 이 한계를
숨기지 않으며, X절의 재현 사례에 포함한다.

\subsection{감사와 사용 차단의 경계}
후속 추적은 도달한 소비 이벤트의 모든 출력을 포함하므로, 문
제 시작점의 디지털 의존 흐름을 차단 대상으로 포착하는 방식
이다. 그러나 실세계 오염 전파의 강도나 제품별 피해 여부는
판정하지 않는다. 문제 시작점의 선택과 차단 승인은 신뢰된 주
체의 판단으로 남는다. 권한 기관의 악의적 행동에 대한 암호학
적 책임 추적도 포함하지 않는다.

동결은 현재 미소비 회계 증표와 미해결 이전 증표의 프로토
콜상 추가 소비를 막는 조치이다. 이미 완료된 소비를 되돌리거
나 발행된 제품 주장을 철회하지 않는다. 물리적 유통 중단을

% Source PDF page 17, column 1
원장 상태 변경만으로 강제하는 것도 아니다. 감사 중 생성된
새 파생 객체는 후속 기준 시점에서 다시 조사해야 하며, 한 실
행이 모든 미래 파생 객체를 원자적으로 봉쇄한다고 주장하지
않는다.

\subsection{구체화해야 할 배포 설정}
정리 13은 명시한 추상 계약과 원장 모델의 결과다. 특히 정리
3은 참여자 비밀·개방값이 공개되지 않는 정직한 이력의 수동
관찰자에 대한 무작위 오라클 모형의 증명이다. 이를 부패 참여
자, 적응적 공개 또는 모든 실제 해시·증명계의 보안으로 확대
하지 않는다. 실제 증명계의 추출 권한·분포 보존과 공동 ZK 모
형, 암호화의 유효 난수 영역 및 키·문맥 계약이 충족되는지 별
도의 구성 대응이 필요하다.

명세의 추상 인터페이스를 구현하려면 다음을 일관되게 고정
해야 한다. 암호 방식과 보안 수준, 곡선·유한체·해시, 정규형과
난수 영역, 수치 단위·배율·비트 수 상한, 교차곱과 합계의 범위,
증명계 초기 설정과 키·공개 입력 순서가 필요하다. 두 키의 역
할은 식 (1)로 정하며, 감사 데이터 암호와 임계 키 보호의 구체
방식·보안 조건, 위원회의 \(K,N\), 인증된 키·암호문 대응과 문맥
인코딩, 초기 비밀 사본 취급 및 외부 키 전달, 키 교체를 지원할
경우의 과거 키 처리, 포함 관계 트리 인코딩, 지원하는 가공 형
상과 정책 상수도 같은 배포 설정에 포함된다.

운영 측면에서는 역할 인증과 적용 범위 자격증명, 블록 높이
기반 에포크 함수와 허용 기한, DPP 필드 추출 규칙, 확정성 및
인증된 과거 상태 조회, 감사 자료 보관과 트랜잭션 확인 한도를
정해야 한다. 이 값을 트랜잭션마다 증명자가 임의로 선택하게
해서는 안 된다. 기존 상태나 암호문의 해석을 바꾸는 변경에는
명시적인 배포·호환 규칙이 필요하다.

\subsection{구현 및 실행 근거의 범위}
명세만으로 회로, 스마트 계약, 클라이언트가 위 조건을 실제로
집행한다고 판정할 수는 없다. 적합성을 확인하려면 아홉 이벤
트의 정상 경로와 함께 잘못된 소유권·문서·입력 중복·회계 범
위·정책·상태·기한의 거부 근거가 필요하다. 생성 때 암호화한
미래 식별자와 실제 소비 식별자가 일치하는지도 실행으로 대
조해야 한다.

감사에서는 분기·합류·중복 도달의 정확한 정점·간선 집합과
미소비 경계 집합, 제품 주장·종료 종단과 기준 시점 격리, 누락
자료의 실패 처리가 검증 대상이다. 동결에서는 실제 추적 결과
가 반환한 전체 대상으로 상태 변경을 실행하고, 부분 실패와 소
비 경합을 유지하며 공통 \(H_{r}\)에서 결과를 확인해야 한다. 차단
된 객체의 후속 소비 거부와 무관한 정상 객체의 소비 허용을 함
께 확인할 필요가 있다. X절은 이 중 선택한 생애주기·추적·완
료 사례를 기호 모델에서 검사한다. 그 결과와 실제 회로·계약·
클라이언트의 적합성 검증은 구분하며, 후자는 본 논문에서 수
행하지 않았다.
